% Source: https://artofproblemsolving.com/community/c3h1642354_geometry
% Stable identity: aops-topic-1642354-q01
\textbf{$EH$ and $FG$.} In triangle $ABD$, $\dfrac{AE}{AB}=\dfrac{AH}{AD}=\dfrac{1}{3}$, so triangles $AEH$ and $ABD$ are similar (SAS) with ratio $\dfrac{1}{3}$. Hence $EH\parallel BD$ and $EH=\dfrac{1}{3}BD$, and $H$ is on the $D$-side of $E$ in the same way as $D$ is on the far side of $B$. Similarly, in triangle $CBD$, $\dfrac{CF}{CB}=\dfrac{CG}{CD}=\dfrac{1}{3}$, so $FG\parallel BD$ and $FG=\dfrac{1}{3}BD$, pointing the same way.

\begin{asy}
usepackage("amsmath");
size(7cm);
defaultpen(fontsize(10pt));
pair A=(0,0),B=(6,-0.6),C=(7,3.6),D=(1.2,4.2);
pair Ee=A+(B-A)/3,F=C+(B-C)/3,G=C+(D-C)/3,H=A+(D-A)/3;
draw(A--B--C--D--cycle,linewidth(0.9));
draw(B--D,gray+dashed);
draw(Ee--F--G--H--cycle,blue+linewidth(0.9));
draw(Ee--G^^F--H,blue+dotted);
pair K=(Ee+G)/2;dot(K);label("$K$",K,1.6*dir(-90));
dot(Ee);dot(F);dot(G);dot(H);
label("$A$",A,SW);label("$B$",B,SE);label("$C$",C,NE);label("$D$",D,NW);
label("$E$",Ee,S);label("$F$",F,E);label("$G$",G,N);label("$H$",H,W);
label("\mbox{}",(min(currentpicture,user=true)+max(currentpicture,user=true))/2,black); // keeps the text after the figure black when LaTeX typesets the labels inline
\end{asy}

In vector form (writing each point for its position vector):
\[
H-E=\left(A+\dfrac{D-A}{3}\right)-\left(A+\dfrac{B-A}{3}\right)=\dfrac{D-B}{3},\qquad
G-F=\left(C+\dfrac{D-C}{3}\right)-\left(C+\dfrac{B-C}{3}\right)=\dfrac{D-B}{3}.
\]
So $\overrightarrow{EH}=\overrightarrow{FG}$: the opposite sides $EH$ and $FG$ of $EFGH$ are equal, parallel and point the same way, and therefore $EFGH$ is a parallelogram.

It is not degenerate: $F-E=\dfrac{2C+B}{3}-\dfrac{2A+B}{3}=\dfrac{2}{3}(C-A)$, so $EF\parallel AC$, while $EH\parallel BD$; the diagonals $AC$ and $BD$ of a convex quadrilateral cross, so they are not parallel, and $E,F,G,H$ are not collinear.

\textbf{The point $K$.} The diagonals of a parallelogram bisect each other, so $K=EG\cap FH$ is the midpoint of both. Indeed $\dfrac{E+G}{2}=\dfrac{F+H}{2}$ follows from $H-E=G-F$. $\blacksquare$
